\section{Introduction}
\IEEEPARstart{I}{ndustrial-agent} deployments must connect each task to its
production context, process signals, and equipment interfaces. Existing work
addresses end-to-end automation control, operator-reviewed control-sequence
generation, modular factory skills, and production scheduling
\cite{xia2025control,huang2026made,koehle2026llmdriven,waseem2026pretrained}.
Other studies target PLC code generation, context-aware control, resource
capability discovery, shop-floor coordination, digital-twin fault handling,
natural-language machine access, and robotic-line control
\cite{liu2026agents4plc,wu2025instructmpc,lim2025llmenabled,zhao2026shopfloor,gill2025llmdt,hofmann2025nlc}.
Across production lines, agent roles, signal maps, protocols, objectives, and
product/run context vary. An integration architecture must accommodate these
differences around shared task and operation services while keeping each
action tied to its active production context. Recent industrial-agent
architectures add AAS-grounded OPC UA actions with signed approval,
P\&ID-grounded proposal checks, and governed sensor--MES--ERP evidence
streams \cite{graul2026aas,schall2026safe,lin2026agentic}. These works address
asset-specific execution, process validation, or closed-loop analytics. This
paper addresses the complementary software-integration question of composing
task and harness routes, production context, semantic protocol I/O, and
route-owned records.

Consider two production scenes with different process signals and interfaces:
cast-film thickness control and annealing-zone temperature control. An agent
task in either scene needs production context, typed signal access, and a
route from a named operation to the configured device protocol. The framework
must also distinguish a command receipt from a subsequent process
observation. The shared task and operation services can therefore be
configured with scene-specific signal maps, protocols, and limits while
retaining a common execution path. The evaluation instantiates this path in
two simulator scenes.

This paper presents AgentWorkShop as a modular framework for task-to-device
integration in industrial agents. It connects collaborative tasks
and registered harnesses to production context, manufacturing execution
system (MES) data, analysis and modeling services, data-acquisition (DAQ) and
write-control (DCW) nodes, and protocol-specific drivers. Its interfaces include
command-line, web, SDK, REST, Model Context Protocol (MCP), and Agent-to-Agent
(A2A) surfaces. Task and chat routes enter the
collaboration runtime, whereas a direct REST tool-invocation route can call
the shared operation dispatcher without creating a task.
Scene blueprints define signal maps, objectives, and guards; channel profiles,
tool policies, and agent bindings scope operations. Agents and callers invoke
named operations through shared runtime checks; the framework does not choose
routes autonomously and retains proposals, command confirmation, and process
observations as distinct records. Fig.~\ref{fig:arch} presents the system
architecture, and Fig.~\ref{fig:runtime} shows the cross-module task flow.
Section~IV details the node-exploration and recipe lifecycles.

This paper asks whether task and direct-operation routes can execute bounded
operations in two simulator scenes with different signal and protocol
mappings while retaining route-specific approval, receipt, and observation
records. It also tests whether approved recipe actions are refused when their
active context changes before software dispatch. A deterministic fixture
separately replays the registered-harness route without external model
decisions. A mock-state suite injects line,
product, recipe, and run changes during approval for \texttt{recipe\_apply},
\texttt{recipe\_trial}, and dispatching \texttt{recipe\_rollback}; all stale
actions are refused before mock-driver writes. A post-recheck fault injection
also tests a context change queued before the dispatch continuation. In a
focused test-only ablation that omits only the tuple-equality check, six
stale product/run changes reach mock actuation, while independent checks still
refuse the other six cases.

The paper makes two contributions:
\begin{enumerate}
\item It specifies and implements a modular task-to-device integration
  contract that
  joins collaborative tasks and registered harnesses to turn-time production
  context, per-call operation checks, semantic I/O, configured protocol
  drivers, and service-owned records. For approved recipe actions, the
  implementation binds approval to a line--product--recipe--run tuple and
  rechecks it after asynchronous authorization validation, immediately before
  entering software dispatch. This is a software-state guard, not an atomic
  lock across device I/O. It
  distinguishes prompt grounding, dispatch authority, device receipt, and
  later process observation; it is an implementation contribution, not a new
  protocol, safety-verification method, or process-control algorithm.
\item It evaluates selected paths in cast-film and annealing simulator scenes.
  In R20, six sequential tasks through one registered harness passed the
  task-to-approval-to-receipt-to-sample linkage checks; none of the six samples
  reached its target band. In R21, a deterministic fixture executed one task
  per scene with no external model and passed 31 runner assertions. R22
  exercised the direct operation route once per scene, with verified
  readbacks and run-tagged snapshots. A separate mock-state suite refused all 12 stale
  line--product--recipe--run mutations on three recipe actuation routes before
  mock-driver writes; a targeted ablation isolates the added coverage for
  product and run changes provided by the final tuple check.
\end{enumerate}

Section~II reviews related industrial-agent and integration work.
Section~III presents the design rationale, architecture, and cross-module
execution flow; Section~IV explains module responsibilities and their
integrated operation. Section~V describes the experiments and results, and
Sections~VI and~VII discuss the evidence limits and conclude the paper.
